% JMLR style article, recreated with standard packages only.
\documentclass[11pt]{article}
\usepackage[letterpaper,textwidth=6in,textheight=8.8in,top=1in]{geometry}
\usepackage{mathptmx}
\usepackage{amsmath}
\usepackage{titlesec}

\pagestyle{plain}

\titleformat{\section}{\large\bfseries}{\thesection.}{0.6em}{}
\titleformat{\subsection}{\normalsize\bfseries}{\thesubsection.}{0.6em}{}
\titlespacing*{\section}{0pt}{1.6ex plus .4ex}{0.9ex}

\begin{document}

\noindent{\small Journal of Machine Learning Research 27 (2026) 1--24
\hfill Submitted 9/25; Revised 4/26; Published 7/26}
\vskip 2pt
\noindent\rule{\textwidth}{0.5pt}

\vskip 22pt
\begin{center}
  {\Large\bfseries Convergence Rates of Stochastic Mirror Descent\\ with Adaptive Regularization\par}
\end{center}

\vskip 14pt
\noindent{\bfseries Nadia Kowalczyk} \hfill {\small\texttt{nadia.kowalczyk@mimuw.edu.pl}}\\
{\small\itshape Institute of Mathematics, University of Warsaw, Warsaw, Poland}
\vskip 8pt
\noindent{\bfseries Julian Weiss} \hfill {\small\texttt{jweiss@ist.ac.at}}\\
{\small\itshape Institute of Science and Technology Austria, Klosterneuburg, Austria}
\vskip 12pt
\noindent{\small\bfseries Editor:} {\small Ravindra Menon}

\vskip 18pt
\begin{center}{\bfseries Abstract}\end{center}
\begin{list}{}{\leftmargin=0.3in\rightmargin=0.3in}\item\relax\small
We analyze stochastic mirror descent when the mirror map itself is adapted
online from observed gradients. Existing analyses fix the geometry in
advance, which forfeits the method's main practical appeal. We show that a
regularizer chosen from a growing family of Bregman divergences achieves the
best rate attainable by any fixed member of the family, up to a logarithmic
factor, without prior knowledge of the gradient distribution. The proof
couples a follow-the-regularized-leader argument with a new stability bound
for time-varying mirror maps. As corollaries we recover AdaGrad-type rates in
the Euclidean case and obtain the first adaptive guarantees for spectral
geometries used in matrix completion.
\end{list}
\vskip 4pt
\noindent{\small{\bfseries Keywords:} mirror descent, adaptive methods, online learning, Bregman divergences, regret bounds}

\section{Introduction}
Mirror descent generalizes gradient descent to non-Euclidean geometries and
underlies most modern adaptive optimizers. Its behavior is governed by the
choice of mirror map, yet that choice is usually made once, before any data
is seen~\cite{ferro2023geometry}. We ask whether the geometry can be learned
during optimization at no asymptotic cost.

With mirror map $\psi_t$ at step $t$, the iterate update is
\begin{equation}
  x_{t+1} \;=\; \arg\min_{x \in \mathcal{X}}
  \Bigl\{ \eta \,\langle \hat{g}_t, x \rangle + D_{\psi_t}(x, x_t) \Bigr\},
  \label{eq:md}
\end{equation}
where $D_{\psi_t}$ is the Bregman divergence of $\psi_t$ and $\hat{g}_t$ is
a stochastic gradient. Our contribution is a schedule for $\psi_t$ under
which Equation~\eqref{eq:md} satisfies an anytime regret bound competitive
with the best fixed geometry in hindsight.

\section{Main Result}
Our central theorem bounds the excess risk after $T$ steps by
$O(\sqrt{\smash[b]{S_T^{\ast} \log T}}/T)$, where $S_T^{\ast}$ is the
cumulative gradient energy measured in the optimal fixed geometry. The
stability argument extends the anytime analysis of~\cite{delgado2024anytime}
and closes a gap left open for spectral regularizers
by~\cite{brekke2022spectral}. All proofs are in Appendix~A, and code
reproducing the matrix-completion experiments accompanies the paper.

\begin{thebibliography}{9}
\bibitem{ferro2023geometry} L. Ferro and D. Anand.
Choosing the geometry of first-order methods.
\textit{Journal of Machine Learning Research}, 24(201):1--52, 2023.
\bibitem{delgado2024anytime} M. Delgado and S. Okabe.
Anytime bounds for follow-the-regularized-leader.
In \textit{Proceedings of the 37th Annual Conference on Learning Theory}, 2024.
\bibitem{brekke2022spectral} A. Brekke and P. Chandrasekar.
Spectral regularization in stochastic approximation.
\textit{Journal of Machine Learning Research}, 23(88):1--39, 2022.
\end{thebibliography}

\end{document}
